\documentclass[aspectratio=169,10pt]{beamer}
% ═══════════════════════════════════════════════════════════════════════════════
% SHARED PREAMBLE — Foundation Models Course (HIT)
% To be \input by each lecture file AFTER \documentclass
% ═══════════════════════════════════════════════════════════════════════════════

% ─────────────────────────────────────────────────────────────────────────────
% BEAMER THEME & BASIC SETUP
% ─────────────────────────────────────────────────────────────────────────────
\usetheme{Madrid}
\usecolortheme{default}

% Remove navigation symbols
\beamertemplatenavigationsymbolsempty

% ─────────────────────────────────────────────────────────────────────────────
% COLORS — HIT Branding
% ─────────────────────────────────────────────────────────────────────────────
\definecolor{hitblue}{RGB}{0,51,102}        % Primary: HIT Navy
\definecolor{hitgold}{RGB}{192,148,0}       % Accent: HIT Gold
\definecolor{hitlight}{RGB}{230,240,250}    % Light background

% Override Beamer palette
\setbeamercolor{primary}{fg=white,bg=hitblue}
\setbeamercolor{secondary}{fg=hitblue,bg=hitlight}
\setbeamercolor{structure}{fg=hitblue}
\setbeamercolor{alerted text}{fg=hitgold}

% ─────────────────────────────────────────────────────────────────────────────
% PACKAGES
% ─────────────────────────────────────────────────────────────────────────────
\usepackage{amsmath}
\usepackage{amssymb}
\usepackage{mathtools}
\usepackage{booktabs}
\usepackage{graphicx}
\usepackage{hyperref}
\usepackage{tikz}
\usepackage{pgfplots}
\usepackage{tcolorbox}
\usepackage{listings}
\usepackage{xcolor}
\usepackage{algorithm}
\usepackage{algpseudocode}

% ─────────────────────────────────────────────────────────────────────────────
% TikZ LIBRARIES
% ─────────────────────────────────────────────────────────────────────────────
\usetikzlibrary{shapes,arrows,positioning,calc,chains,scopes}
\pgfplotsset{compat=1.17}

% ─────────────────────────────────────────────────────────────────────────────
% CODE LISTINGS
% ─────────────────────────────────────────────────────────────────────────────
\lstset{
  basicstyle=\ttfamily\small,
  breaklines=true,
  commentstyle=\color{gray},
  keywordstyle=\color{hitblue},
  stringstyle=\color{hitgold},
  showstringspaces=false,
  backgroundcolor=\color{hitlight},
  frame=single,
  rulecolor=\color{hitblue}
}

% ─────────────────────────────────────────────────────────────────────────────
% TCOLORBOX STYLES
% ─────────────────────────────────────────────────────────────────────────────
\tcbset{
  hitbox/.style={
    colback=hitlight,
    colframe=hitblue,
    fonttitle=\bfseries,
    boxrule=1pt
  },
  alertbox/.style={
    colback=yellow!10,
    colframe=hitgold,
    fonttitle=\bfseries,
    boxrule=1.5pt
  }
}

% ─────────────────────────────────────────────────────────────────────────────
% CUSTOM COMMANDS
% ─────────────────────────────────────────────────────────────────────────────

% Placeholder for figures
\newcommand{\placeholder}[3]{
  \begin{center}
    \fbox{\begin{minipage}[c][#2][c]{#1}
      \centering\small\color{gray}[Figure: #3]
    \end{minipage}}
  \end{center}
}

% Section divider frame
\newcommand{\sectionframe}[2]{
  \begin{frame}[plain]
    \begin{tikzpicture}[remember picture,overlay]
      \fill[hitblue] (current page.south west) rectangle (current page.north east);
      \node[anchor=center, white, font=\Large\bfseries] at (current page.center) {
        \begin{tabular}{c}
          #1 \\[0.3cm]
          {\small\color{hitgold}#2}
        \end{tabular}
      };
    \end{tikzpicture}
  \end{frame}
}

% ─────────────────────────────────────────────────────────────────────────────
% LOAD CUSTOM MACROS
% ─────────────────────────────────────────────────────────────────────────────
% ═══════════════════════════════════════════════════════════════════════════════
% SHARED MACROS — Mathematical Notation for Foundation Models Course
% ═══════════════════════════════════════════════════════════════════════════════

% ─────────────────────────────────────────────────────────────────────────────
% ACTIVATION & LAYER FUNCTIONS
% ─────────────────────────────────────────────────────────────────────────────
\newcommand{\softmax}{\mathrm{softmax}}
\newcommand{\sigmoid}{\mathrm{sigmoid}}
\newcommand{\relu}{\mathrm{ReLU}}
\newcommand{\gelu}{\mathrm{GELU}}
\newcommand{\LayerNorm}{\mathrm{LayerNorm}}
\newcommand{\FFN}{\mathrm{FFN}}
\newcommand{\MHA}{\mathrm{MHA}}
\newcommand{\Attn}{\mathrm{Attention}}

% ─────────────────────────────────────────────────────────────────────────────
% ATTENTION COMPONENTS
% ─────────────────────────────────────────────────────────────────────────────
\newcommand{\query}{Q}
\newcommand{\key}{K}
\newcommand{\val}{V}
\newcommand{\dmodel}{d_{\text{model}}}
\newcommand{\dk}{d_k}
\newcommand{\nheads}{h}

% Scaled dot-product attention formula
\newcommand{\SDPA}{
  \Attn(\query, \key, \val) = \softmax\left(\frac{\query \key^T}{\sqrt{\dk}}\right) \val
}

\newcommand{\CrossAttn}{\text{CrossAttention}}

% ─────────────────────────────────────────────────────────────────────────────
% PROBABILITY & EXPECTATION
% ─────────────────────────────────────────────────────────────────────────────
\newcommand{\E}{\mathbb{E}}
\newcommand{\Prob}{\mathbb{P}}
\newcommand{\KL}{\mathrm{KL}}
\newcommand{\ELBO}{\mathrm{ELBO}}
\newcommand{\given}{\mid}

% ─────────────────────────────────────────────────────────────────────────────
% NORMS & OPERATIONS
% ─────────────────────────────────────────────────────────────────────────────
\newcommand{\norm}[1]{\left\lVert #1 \right\rVert}
\newcommand{\abs}[1]{\left| #1 \right|}
\newcommand{\inner}[2]{\langle #1, #2 \rangle}
\newcommand{\T}{^{\top}}

% ─────────────────────────────────────────────────────────────────────────────
% VECTORS & MATRICES
% ─────────────────────────────────────────────────────────────────────────────
\newcommand{\bvec}[1]{\mathbf{#1}}
\newcommand{\mat}[1]{\mathbf{#1}}
\newcommand{\iid}{\stackrel{\text{i.i.d.}}{\sim}}
\newcommand{\defeq}{\coloneq}

% ─────────────────────────────────────────────────────────────────────────────
% LANGUAGE MODELING
% ─────────────────────────────────────────────────────────────────────────────
\newcommand{\vocab}{\mathcal{V}}
\newcommand{\context}{\text{context}}
\newcommand{\token}{t}
\newcommand{\seq}{x}
\newcommand{\ptheta}{p_\theta}
\newcommand{\pref}{p_{\text{ref}}}

% ─────────────────────────────────────────────────────────────────────────────
% RLHF & DPO
% ─────────────────────────────────────────────────────────────────────────────
\newcommand{\piref}{\pi_{\text{ref}}}
\newcommand{\pitheta}{\pi_\theta}
\newcommand{\yw}{y_w}
\newcommand{\yl}{y_l}
\newcommand{\DPOloss}{\mathcal{L}_{\mathrm{DPO}}}

% ─────────────────────────────────────────────────────────────────────────────
% RETRIEVAL & RAG
% ─────────────────────────────────────────────────────────────────────────────
\newcommand{\docs}{\mathcal{D}}
\newcommand{\qvec}{q}
\newcommand{\dvec}{d}

% ─────────────────────────────────────────────────────────────────────────────
% AGENTS & REASONING
% ─────────────────────────────────────────────────────────────────────────────
\newcommand{\obs}{o}
\newcommand{\act}{a}
\newcommand{\thought}{t}
\newcommand{\tools}{\mathcal{T}}
\newcommand{\history}{h}
\newcommand{\concat}{\oplus}

% ─────────────────────────────────────────────────────────────────────────────
% SETS & LOGIC
% ─────────────────────────────────────────────────────────────────────────────
\newcommand{\R}{\mathbb{R}}
\newcommand{\N}{\mathbb{N}}
\newcommand{\Z}{\mathbb{Z}}

\endinput


% ─────────────────────────────────────────────────────────────────────────────
% TITLE & AUTHOR (can be overridden in individual lectures)
% ─────────────────────────────────────────────────────────────────────────────
\author[D. Lederman]{Dror Lederman}
\institute{Holon Institute of Technology (HIT) \\ Data Science Department}
\date{\today}

% ─────────────────────────────────────────────────────────────────────────────
% FOOTER & HEADER
% ─────────────────────────────────────────────────────────────────────────────
\setbeamertemplate{footline}{
  \leavevmode%
  \hbox{%
    \begin{beamercolorbox}[wd=0.33\paperwidth,ht=2.25ex,dp=1ex,center]{author in head/foot}%
      \usebeamerfont{author in head/foot}\insertshortauthor
    \end{beamercolorbox}%
    \begin{beamercolorbox}[wd=0.34\paperwidth,ht=2.25ex,dp=1ex,center]{title in head/foot}%
      \usebeamerfont{title in head/foot}\insertshorttitle
    \end{beamercolorbox}%
    \begin{beamercolorbox}[wd=0.33\paperwidth,ht=2.25ex,dp=1ex,right]{frame number}%
      \usebeamerfont{frame number}\insertframenumber{} / \inserttotalframenumber\hspace*{2ex}
    \end{beamercolorbox}%
  }%
  \vskip0pt%
}

\endinput


\title{Week 4: Multimodal Foundation Models}
\subtitle{Foundation Models and Agentic AI}
\author{Lecturer Name}
\institute{Holon Institute of Technology (HIT)}
\date{Semester, 2025}

\begin{document}

\begin{frame}
  \titlepage
\end{frame}

\begin{frame}{Learning Objectives}
  After this lecture you will be able to:
  \begin{itemize}
    \item Explain \textbf{vision-language alignment} and how CLIP achieves it.
    \item Describe the \textbf{Flamingo} architecture and its cross-attention mechanism.
    \item Compare \textbf{late-fusion} and \textbf{early-fusion} multimodal architectures.
    \item Understand the design of \textbf{GPT-4o / Gemini-style} unified multimodal models.
    \item Identify key challenges in \textbf{multimodal reasoning} and evaluation.
  \end{itemize}
\end{frame}

\section{Vision-Language Models: Foundations}
\sectionframe{Vision-Language Alignment}{Teaching models to connect images and text}

\begin{frame}{Why Multimodal?}
  \begin{columns}[T]
    \column{0.52\textwidth}
      \textbf{Motivation:}
      \begin{itemize}
        \item Human intelligence is fundamentally multimodal.
        \item Language alone is ambiguous: ``the bank'' (river vs.\ financial).
        \item Many real tasks require visual grounding: medical imaging, robotics, UI automation.
        \item Image captions, VQA, document understanding.
      \end{itemize}
      \vspace{0.4em}
      \textbf{Historical approaches:}
      \begin{itemize}
        \item Feature concatenation (CNN + LSTM): ViLBERT, LXMERT
        \item Limited by separately pretrained unimodal encoders
      \end{itemize}
    \column{0.44\textwidth}
      \placeholder{0.95\textwidth}{4cm}{Multimodal task examples: VQA, captioning, OCR}
  \end{columns}
\end{frame}

\begin{frame}{Vision Encoders}
  \begin{tcolorbox}[hitbox, title=Vision Transformer (ViT)]
    Divide image into $N$ non-overlapping patches of size $P \times P$.
    Linearly project each patch; add position embeddings; feed into a transformer encoder.
    \[
      \bvec{z}_0 = [\bvec{x}_\text{cls}; \bvec{x}_1 \mat{E}; \ldots; \bvec{x}_N \mat{E}] + \mat{E}_\text{pos}
    \]
  \end{tcolorbox}
  \vspace{0.4em}
  \begin{itemize}
    \item Patch size $P=16$ (ViT-B/16): $224^2 / 16^2 = 196$ patches + 1 CLS token.
    \item No convolutional inductive bias — learns spatial structure from data.
    \item Pre-trained ViT features are the backbone of CLIP, Flamingo, LLaVA, etc.
  \end{itemize}
\end{frame}

\section{CLIP: Contrastive Language-Image Pretraining}
\sectionframe{CLIP}{Aligning vision and language with contrastive learning}

\begin{frame}{CLIP Architecture}
  \textbf{Radford et al.\ (2021)}~\cite{radford2021learning} — OpenAI.
  \vspace{0.4em}
  \begin{columns}[T]
    \column{0.55\textwidth}
      \textbf{Setup:}
      \begin{itemize}
        \item Dual encoder: image encoder (ViT or ResNet) + text encoder (Transformer).
        \item Train on 400M image-caption pairs from the internet.
        \item Maximize cosine similarity of matched pairs; minimize for unmatched.
      \end{itemize}
      \vspace{0.4em}
      \textbf{Loss} (InfoNCE / NT-Xent):
      \[
        \mathcal{L} = -\frac{1}{N}\sum_{i=1}^N \log
        \frac{\exp(\bvec{i}_i \cdot \bvec{t}_i / \tau)}
             {\sum_{j=1}^N \exp(\bvec{i}_i \cdot \bvec{t}_j / \tau)}
      \]
    \column{0.42\textwidth}
      \placeholder{0.95\textwidth}{4cm}{CLIP: contrastive matrix of image-text similarities}
  \end{columns}
  \note{$\tau$ is a learned temperature. Low $\tau$ → peaked distribution; high $\tau$ → uniform.}
\end{frame}

\begin{frame}{CLIP: Zero-Shot Transfer}
  \begin{tcolorbox}[hitbox, title=Zero-Shot Classification]
    For each class label $c$, form prompt: \textit{``A photo of a \{c\}.''}\\
    Encode all prompts. Classify image by nearest-neighbor in embedding space.
  \end{tcolorbox}
  \vspace{0.4em}
  \textbf{Results:}
  \begin{itemize}
    \item Zero-shot CLIP matches supervised ResNet-50 on ImageNet.
    \item Transfers to 30+ datasets without any task-specific training.
    \item Demonstrates that \textbf{natural language is a rich supervision signal}.
  \end{itemize}
  \vspace{0.3em}
  \textbf{Limitations:}
  \begin{itemize}
    \item Does not understand compositionality well: ``red square above blue circle'' fails.
    \item Biases from internet-scale image-text data.
  \end{itemize}
\end{frame}

\section{Flamingo-Style Models}
\sectionframe{Flamingo}{Few-shot multimodal learning via cross-attention}

\begin{frame}{Flamingo Architecture}
  \textbf{Alayrac et al.\ (2022)}~\cite{alayrac2022flamingo} — DeepMind.
  \vspace{0.4em}
  \begin{columns}[T]
    \column{0.55\textwidth}
      \textbf{Design:}
      \begin{itemize}
        \item Frozen large LLM backbone (70B Chinchilla).
        \item Frozen CLIP-style vision encoder.
        \item \textbf{Perceiver Resampler}: compresses variable-length image features to fixed-length visual tokens.
        \item \textbf{Cross-attention layers} interleaved with LLM layers — only these are trained.
      \end{itemize}
      \vspace{0.3em}
      \textbf{Key insight:} Freeze unimodal experts; train only the bridge.
    \column{0.42\textwidth}
      \placeholder{0.95\textwidth}{4cm}{Flamingo: frozen LLM + cross-attention to visual tokens}
  \end{columns}
  \note{The Perceiver Resampler is a compact cross-attention module that produces a fixed number of output tokens regardless of image resolution.}
\end{frame}

\begin{frame}{In-Context Multimodal Learning}
  Flamingo supports \textbf{few-shot multimodal prompting}:
  \begin{tcolorbox}[hitbox]
    \texttt{[Image: cat] Caption: A fluffy orange cat.}\\
    \texttt{[Image: dog] Caption: A brown labrador.}\\
    \texttt{[Image: \textit{test image}] Caption:}$\square$
  \end{tcolorbox}
  \vspace{0.4em}
  \begin{itemize}
    \item Each image is treated as a special token sequence in the LLM context.
    \item The model learns to condition generation on visual input from in-context examples.
    \item Flamingo-80B achieves state-of-the-art on VQA, captioning with 4 shots.
  \end{itemize}
\end{frame}

\section{Modern Unified Multimodal Models}
\sectionframe{Unified Multimodal Models}{GPT-4V, Gemini, LLaVA}

\begin{frame}{LLaVA: Visual Instruction Tuning}
  \textbf{Liu et al.\ (2023)}~\cite{liu2023llava}
  \vspace{0.4em}
  \begin{columns}[T]
    \column{0.55\textwidth}
      \textbf{Architecture:}
      \begin{itemize}
        \item CLIP ViT-L/14 → linear projection → LLaMA token space.
        \item Instruction-tune the combined model on GPT-4 generated visual instructions.
        \item Open-source, reproducible, 7B parameters.
      \end{itemize}
      \vspace{0.4em}
      \textbf{Key insight:} High-quality instruction data matters more than architecture.
    \column{0.42\textwidth}
      \placeholder{0.95\textwidth}{3.5cm}{LLaVA: CLIP features → projection → LLaMA}
  \end{columns}
\end{frame}

\begin{frame}{GPT-4o and Gemini: Native Multimodality}
  \begin{columns}[T]
    \column{0.52\textwidth}
      \textbf{GPT-4 Technical Report}~\cite{openai2023gpt4}:
      \begin{itemize}
        \item Accepts image and text interleaved.
        \item Details of architecture not disclosed.
        \item Strong on chart/diagram understanding, OCR, spatial reasoning.
      \end{itemize}
      \vspace{0.4em}
      \textbf{Gemini 1.5 / GPT-4o:}
      \begin{itemize}
        \item Process audio, video, images, and text natively.
        \item Up to 1M token context window (Gemini 1.5).
        \item Unified encoder: no separate vision/text pathways.
      \end{itemize}
    \column{0.44\textwidth}
      \placeholder{0.95\textwidth}{4cm}{GPT-4o: unified any-to-any modality model}
  \end{columns}
  \note{The trend is toward ``any-to-any'' models that process and generate across all modalities. This is an active research area.}
\end{frame}

\begin{frame}{Multimodal Reasoning Challenges}
  \begin{tcolorbox}[alertbox, title=Current Limitations]
    \begin{itemize}
      \item \textbf{Compositionality}: struggling with spatial relations, counting, multi-step visual reasoning.
      \item \textbf{Hallucination}: confidently describing objects not in the image.
      \item \textbf{Fine-grained OCR}: reading dense text in images is still error-prone.
      \item \textbf{Video understanding}: temporal reasoning across frames is nascent.
      \item \textbf{3D spatial understanding}: depth, perspective, rotation.
    \end{itemize}
  \end{tcolorbox}
\end{frame}

\section{Paper Discussion}
\begin{frame}{Paper Discussion}
  \papercite{Learning Transferable Visual Models From Natural Language Supervision (CLIP)}
    {Radford, Kim, Hallacy, Ramesh, Goh et al.\ (OpenAI)}
    {ICML 2021~\cite{radford2021learning}}
  \vspace{0.4em}
  \textbf{What made CLIP revolutionary?}
  \begin{itemize}
    \item Showed that internet-scale noisy supervision is sufficient for strong visual representations.
    \item Zero-shot transfer without task-specific heads.
    \item The contrastive objective became the foundation of all subsequent VLMs.
  \end{itemize}
  \vspace{0.3em}
  \textbf{Discussion:} What are the societal implications of training on billions of internet images?
  Who controls the ``visual world model'' encoded in CLIP?
\end{frame}

\section{Lab / Demo}
\begin{frame}{Lab Idea: Multimodal Probing with LLaVA}
  \begin{tcolorbox}[hitbox, title=Lab 4 — Multimodal VQA]
    \textbf{Goal:} Probe the multimodal capabilities and failure modes of an open-source VLM.\\[4pt]
    \textbf{Tools:} \texttt{LLaVA-1.5-7B} via HuggingFace or Ollama
    \begin{enumerate}
      \item Test basic VQA: object identification, color, counting.
      \item Test spatial reasoning: ``Is the cup to the left or right of the book?''
      \item Test compositional queries: ``How many red circles are above the blue square?''
      \item Attempt to elicit hallucinations: ask about objects not in the image.
      \item Document: where does the model succeed/fail, and why?
    \end{enumerate}
    \textbf{Deliverable:} Annotated test set (20 examples) + failure analysis.
  \end{tcolorbox}
\end{frame}

\section{Discussion \& Summary}
\begin{frame}{Discussion Questions}
  \begin{enumerate}
    \item CLIP uses contrastive learning on image-caption pairs. What are the data biases inherent in web-crawled image-caption data? How do they manifest downstream?
    \item Flamingo freezes the LLM backbone and only trains cross-attention layers. What are the advantages and disadvantages of this ``frozen backbone'' approach?
    \item GPT-4o can process video, audio, and text simultaneously. What new capabilities does this enable? What new risks does it create?
    \item Multimodal models can hallucinate objects not present in an image. Is this the same phenomenon as text hallucination, or something different?
  \end{enumerate}
\end{frame}

\begin{frame}{Summary}
  \begin{itemize}
    \item \textbf{Vision encoders} (ViT) provide patch-level features; \textbf{CLIP} aligns them with text.
    \item \textbf{CLIP's contrastive loss} enables zero-shot image classification via text prompts.
    \item \textbf{Flamingo} bridges frozen unimodal experts with cross-attention + Perceiver Resampler.
    \item \textbf{LLaVA} shows instruction tuning with quality data achieves strong open-source VLM.
    \item Modern models (GPT-4o, Gemini) aim for native, unified any-to-any multimodality.
    \item Key challenges: compositionality, hallucination, video reasoning.
  \end{itemize}
  \vspace{0.3em}
  \textbf{Next week:} Embeddings and semantic search — the retrieval infrastructure behind RAG.
\end{frame}

\begin{frame}[allowframebreaks]{Suggested Reading}
  \nocite{radford2021learning, alayrac2022flamingo, liu2023llava, openai2023gpt4}
  \bibliography{../../references}
\end{frame}

\end{document}
